Al laboratori dissenyem un experiment amb dues masses que fan moviments independents. La primera massa és de 0,5~kg i penja d'una molla vertical. La segona massa es troba subjectada a un disc vertical que gira a una velocitat angular de 6,41~rad/s i el centre d'aquesta massa és a una distància de 19~cm del centre del disc. Ambdues masses s'il·luminen lateralment i s'observa que les seves ombres segueixen exactament el mateix moviment harmònic simple. Per a aconseguir-ho, deixem anar des de baix la massa de la molla just a $t = 0\ \mathrm{s}$, moment en què la massa del disc també passa pel punt més baix de la rotació.

\figura[0.45]{fig1.png}

\begin{apartat}{1,25}
Escriviu l'equació de la posició vertical de les ombres respecte al temps. Trobeu la constant elàstica de la molla i l'energia mecànica del moviment harmònic simple.
\end{apartat}

\begin{apartat}{1,25}
Calculeu l'equació de la velocitat i l'energia cinètica respecte al temps de la massa que penja de la molla. Representeu en la quadrícula adjunta l'energia mecànica, l'energia potencial i l'energia cinètica en funció de la posició vertical per a la massa que penja de la molla.
\end{apartat}

\figura[0.75]{fig2.png}
